\documentclass[11pt]{article}
\usepackage[letterpaper,margin=.68in,headheight=16pt,headsep=.18in,footskip=.28in]{geometry}
\usepackage{amsmath,array,enumitem,fancyhdr,hyperref}
\pagestyle{fancy}\fancyhf{}\renewcommand{\headrulewidth}{.3pt}
\fancyhead[L]{\small Bellingham Math Circle / Week 13 / Return visits}
\fancyhead[R]{\small Adult guide}
\fancyfoot[L]{\small RV-13-FAC-v1 / Draft and unpiloted}\fancyfoot[R]{\small\thepage}
\setlength{\emergencystretch}{2em}
\setlength{\parindent}{0pt}\setlength{\parskip}{7pt}
\setlist[itemize]{leftmargin=1.3em,itemsep=3pt,topsep=4pt}
\hypersetup{hidelinks,pdftitle={Week 13 return-visit facilitator guide}}
\begin{document}
{\Large\bfseries Mathematical overview}\par\medskip
\textbf{Facts and limits.} A packing of edge-disjoint routes can share internal dots. If internal dots instead have capacity one, the maximum can drop. The central-hub board has two edge-disjoint routes and only one internally vertex-disjoint route; adding the bypass A\(\to\)C permits two of the latter. Edge cuts and vertex cuts answer different questions.

Two routes with specified endpoints need not fit even when the same network carries two routes after reassigning the endpoints. On the P,Q,X,Y board, P\(\to\)Y and Q\(\to\)X both require the single arrow C\(\to\)D, whereas P\(\to\)X and Q\(\to\)Y fit. The single-start max-flow theorem does not promise simultaneous routes for prescribed pairs.

On a finite directed network with one start and one finish and nonnegative integer arrow capacities, an arrow holds that many simultaneous routes. Replacing capacity \(c\) by \(c\) separate lanes gives integral packing; the maximum equals the minimum sum of cut capacities. The printed capacity board supports four routes, and increasing C\(\to\)T from three to four permits five. Length and sequential travel do not represent capacity here.

\textbf{Children's destination.} Identify exactly what is occupied, distinguish endpoint assignment from route choice, and prove an optimum by exhibiting routes and a matching obstruction.

\textbf{Entry map.} Problem 1: K--1 and up with oral rules and occupied-dot counters. Problem 2: grades 2--3 and up, preserve both endpoint obligations simultaneously. Problem 3: grades 4--5, small totals, lane counts and cut certificates. No flow algebra is required.

\textbf{New versus base.} The base uses one start/finish and unit edge capacity, including rerouting and cuts. These visits alter vertex sharing, impose two endpoint pairs, and introduce integer lanes.
\newpage {\Large\bfseries Prepare a return visit}\par\medskip
\textbf{Status and use.} Three investigations extend one familiar theme. These companions are separate from the base packets and may support several visits. All pages are new and unpiloted. Printed labels describe prerequisites, not age restrictions. Give one page at a time; completing every page is not the goal.

\textbf{Materials and preparation.} For each pair, pencil/eraser, two route colors or two clear line styles, and 16 markers no wider than 15 mm. Reserve six boards/marker sets for eleven children. For vertex exclusion use a marker on every occupied internal dot; for capacities place one short tick or token in each reserved lane slot. Provide spare paper. Keep arrowheads visible and distinguish crossings without dots from junctions.

\textbf{Staffing.} Current context: eleven children, fixed KK11 / 3333 / 445 tables, one adult anchored to each table. Use pairs; the upper third child rotates as reader or checker. Routine adult reading, arrow drawing or recording can leave the mathematical decisions with children. If adults must choose every move, use the earlier entry rather than run the investigation for them.

\textbf{Short common launch.} Let pairs trace one route. Reserve its whole path before tracing another. Demonstrate an unused arrow entering an already occupied dot: under edge-sharing rules this can be legal; under the new dot rule it is blocked. Put a marker on that dot so a partner enforces the rule. On a spare arrow with two lane boxes, put one reservation in each and show that a third cannot fit.

\textbf{Suggested first route.} Begin with the hub board and physical dot reservations. The two-pair page needs a child who can retain the pair obligations while rerouting; the reader or adult may restate them. Offer capacity bookkeeping only after simultaneous reservation is secure.

\textbf{Flexible timing.} Allow 3--5 minutes of material play and 2--4 minutes for the common demonstration, then 15--25 minutes for one investigation and 5 minutes to share. A longer second visit can spend 30--45 minutes on the same investigation, including unfinished examples or its explanation. These are planning estimates, not observed timings. Do not require all three within one hour.

\textbf{Questions and hints.} Start with ``What can you try?'' and ``What would count as success?'' Check the actual rule before giving a strategy. Optional questions and hints are keyed below; use them only after a real attempt. A counter argument or drawing may be a complete explanation.

\textbf{Source and pilot record.} Mathematical examples and arguments here are original local follow-ups to the current base theme. Consulted teaching source: \emph{Math Circle by the Bay}, Preface pp. vii--x (local PDF pp. 8--11), on interaction, manipulatives, hints, explanations and themes spanning multiple years. Our exact launch, entry route and timing are design inferences. Year-1 \emph{Handouts 6}, Problems 6.2--6.6, provided a local example of connecting representations. Record actual problem, starting route, examples tried, what needed adult rescue, and what children want to resume. Distinguish observations from hypotheses. Counter fit, materials stock, procedures and classroom response have not been physically tested.
\newpage
\textbf{Problem 1: roads versus junctions.} On the hub board use S-A-H-C-T and S-B-H-D-T. They share H but no arrow, so two edge-disjoint routes fit. Three are impossible because only two arrows leave S.

When no internal dot may be shared, only one fits: every S-to-T route contains H. Closing H stops all routes, giving the matching one-dot obstruction. Closing one arrow cannot do the same on this board, since the exhibited two routes have no common arrow.

Add A\(\to\)C. Now S-A-C-T and S-B-H-D-T share no internal dot and fit together. Again two is maximal because each route needs a distinct first arrow from S.

\textbf{If needed.} Put a counter on H after the first route. Ask whether the second route needs that exact dot. Do not model these as pawns taking successive journeys: successive routes could reuse both roads and dots and would change the question.

\textbf{Extension.} Invent a board where deleting one internal dot stops all travel but at least three arrows must be closed. A three-branch hub supplies an example; the child should demonstrate three disjoint arrow routes and show that every route meets the hub.

\textbf{Limits.} Internally vertex-disjoint routes may share S and T. If a direct S-to-T edge exists, removing internal vertices alone cannot block it; a general vertex-cut formulation must handle that case separately.
\newpage
\textbf{Problem 2: two assigned deliveries.} The direct assignments fit: P-A-X and Q-B-Y use no common arrow. Each crossed assignment exists alone: P-A-C-D-Y and Q-B-C-D-X.

They cannot both be reserved simultaneously. A crossed route leaving A must go to C, because A-X reaches the wrong terminal and terminates there. A crossed route leaving B likewise must go to C. Both then must use C-D. This arrow has capacity one, so the pair cannot fit. This checks all possible crossed routes, rather than merely rejecting the drawn attempts.

If the endpoints can be reassigned, two routes fit again. There is no contradiction with max-flow/min-cut: that theorem joins a single source to a single sink (or handles unassigned aggregate supplies), not specified source-to-terminal pairings.

\textbf{If needed.} Keep a card P\(\to\)Y beside one marker and Q\(\to\)X beside the other. Ask a child to check the finish before reserving the route. Never update which terminal is required during a run.

\textbf{Extension.} What one-arrow addition permits both crossed deliveries? A\(\to\)Y suffices: P-A-Y and Q-B-C-D-X. B\(\to\)X works symmetrically. The new arrow is a changed network, not an unmarked jump allowed on the original.

\textbf{Observe.} Routine adult reading is helpful; if the adult must choose and reroute both paths, delay the simultaneous-pair investigation and return to the first page.
\newpage
\textbf{Problem 3: several lanes on one road.} Capacities are S-A:3, S-B:2, A-C:2, B-C:2, C-T:3, A-T:1. One route S-A-T, two copies of S-A-C-T and one S-B-C-T give four. Copies are separate units using distinct available lanes.

Only four units can enter T, through capacities three and one. This cut proves that five cannot fit and certifies that the displayed four are optimal. A maximum collection could distribute the three C-T units differently; it need not use every outgoing lane of S.

Increasing C-T by one permits five: one S-A-T, two S-A-C-T and two S-B-C-T. All lane capacities are respected. The old T-entry cut now has capacity five; the S-exit cut also has five, proving optimality. Zero changes cannot suffice, so one lane increase is minimal.

\textbf{If needed.} Draw distinct lane slots beside an arrow. A partner crosses out a slot when a whole route reserves it. Keep just a route list and the slots, not several simultaneous cumulative tables.

\textbf{General continuation.} Replace every c-lane arrow by c separate parallel one-lane arrows. The base unit-capacity theorem then applies; minimum cut size becomes the sum of the original capacities. This argument requires nonnegative integer capacities and one start and one finish. It does not repair the prescribed-pair obstruction from Problem 2.

\textbf{Checks.} The source audit independently enumerates simple routes and their multiplicities, plus all start-side cuts. It confirms capacities four and five and the dot/endpoint obstructions.
\end{document}
